\documentclass[11pt,a4paper]{article}

\usepackage[margin=2.5cm]{geometry}
\usepackage{amsmath,amssymb,amsthm}
\usepackage{bm}
\usepackage{booktabs}
\usepackage{xcolor}
\usepackage{hyperref}
\usepackage{enumitem}

\newtheorem{remark}{Remark}
\newtheorem{assumption}{Assumption}

\newcommand{\R}{\mathbb{R}}
\newcommand{\dom}{\Omega}
\newcommand{\pore}[1]{\Omega_{#1}}
\newcommand{\iface}{\gamma}
\newcommand{\vel}{\bm{u}}
\newcommand{\pres}{p}
\newcommand{\stress}{\bm{\sigma}}
\newcommand{\normal}{\bm{n}}
\newcommand{\trac}{\bm{\lambda}}
\newcommand{\DtN}{\mathbf{G}}
\newcommand{\Schur}{\mathbf{S}}
\newcommand{\one}{\mathbf{1}}
\newcommand{\zero}{\bm{0}}
\newcommand{\modeset}{\mathcal{M}}
\newcommand{\loadvec}{\bm{f}}
\newcommand{\rank}{\operatorname{rank}}

\hypersetup{colorlinks=true,linkcolor=blue,urlcolor=blue,citecolor=blue}

\begin{document}

\title{\textbf{Affine-DDPNM: Methodology}\\[4pt]
       \large Version 6 --- revised 2026-08-07}
\author{}
\date{}
\maketitle

\begin{center}
\fbox{\begin{minipage}{0.92\textwidth}
This revision fixes the traction/load sign convention, the orientation of
tangential modes, the assumptions for local invertibility and global SPD,
and the interpretation of the Galerkin enrichment statement.
\end{minipage}}
\end{center}

We extend the domain-decomposition pore-network method (DD-PNM) of
\cite{sun2025ddpnm} by enriching the interface traction space from a single
constant-normal mode to a nine-mode affine representation.

\begin{center}
\fbox{\begin{minipage}{0.92\textwidth}
\textbf{Scope.}  The formulation and the accompanying Galerkin analysis
are stated for \textbf{planar interfaces} with a fixed orthonormal frame
and dimensionless coordinates $(s,t)$.  The watershed implementation on
non-planar, multi-facet interfaces is a numerical extension and is
\emph{not} covered by the present approximation theorem.
\end{minipage}}
\end{center}
\smallskip

% ======================================================================
\section{Model problem and weak-form convention}
% ======================================================================

Steady incompressible Stokes flow in $\dom \subset \R^3$:
\begin{align}
  -\nu\Delta\vel + \nabla\pres &= \zero,                && \text{in } \dom,\\
  \nabla\cdot\vel &= 0,                                   && \text{in } \dom,\\
  \vel &= \zero,                                          && \text{on } \Gamma_{\text{wall}},\\
  \stress(\vel,\pres)\,\normal &= -p_b\,\normal,         && \text{on } \Gamma_{\text{in}}\cup\Gamma_{\text{out}},
\end{align}
where $\stress(\vel,\pres) := -\pres I + \nu\nabla\vel$ is the pseudostress.
The inlet/outlet condition prescribes the normal pseudotraction, not the
pressure directly: $\normal\cdot\stress\normal = -p_b$ does not imply
$\pres = p_b$ pointwise because the normal viscous term contributes.
The incompressible weak formulation uses
$b_i(\bm{v},q) = -\int_{\pore{i}} q\,\nabla\cdot\bm{v}\,\mathrm{d}x$,
giving the local discrete system
\begin{equation}\label{eq:local-fe}
  \begin{bmatrix} A_i & B_i^\top \\ B_i & 0 \end{bmatrix}
  \begin{bmatrix} U_i \\ P_i \end{bmatrix}
  = \begin{bmatrix} f_i \\ \zero \end{bmatrix}.
\end{equation}
No pressure stabilisation is used.

% ======================================================================
\section{Domain partition and oriented interface frames}
% ======================================================================

$\dom$ is partitioned into $N$ non-overlapping subdomains
$\{\pore{i}\}_{i=1}^N$.  Each internal interface $\iface_k$ is planar
and shared by pores $(i_k,j_k)$.  Choose once a global orthonormal frame
$\{\normal_k,\bm{t}_{k,1},\bm{t}_{k,2}\}$ with $\normal_k$ directed from
$i_k$ to $j_k$.  Let $\sigma_{i,k}=+1$ for $i=i_k$ and $-1$ for $i=j_k$.

To match the code's assembled normal direction $-\normal_{i,k}$ and
tangential directions $\sigma_{i,k}\bm{t}_{k,\nu}$, define
\textbf{oriented inward-traction directions}
\begin{equation}\label{eq:dir}
  \bm{d}_{i,k,n} = \normal_{i,k},\qquad
  \bm{d}_{i,k,t_\nu} = -\sigma_{i,k}\,\bm{t}_{k,\nu}\;\;(\nu=1,2).
\end{equation}
The minus sign in the tangential definition is an orientation choice;
changing it together with the corresponding coefficient changes no
traction space.  Crucially,
\begin{equation}\label{eq:opposite}
  \bm{d}_{j,k,\beta} = -\bm{d}_{i,k,\beta}
  \qquad (\beta = n,t_1,t_2).
\end{equation}

Half-span--normalised in-plane coordinates
$s_k(\bm{x}) = ((\bm{x}-\bm{c}_k)\cdot\bm{t}_{k,1})/a_k$,
$t_k(\bm{x}) = ((\bm{x}-\bm{c}_k)\cdot\bm{t}_{k,2})/b_k$,
with scalar shapes $\varphi_0=1$, $\varphi_s=s_k$, $\varphi_t=t_k$,
$\modeset = \{0,s,t\}\times\{n,t_1,t_2\}$.

\textbf{Column ordering.}
The \emph{code} uses component-major ordering
\[
\begin{aligned}
\mathcal O_{\mathrm{code}}={}&
\bigl((0,n),(s,n),(t,n),\\
& \;(0,t_1),(s,t_1),(t,t_1),\\
& \;(0,t_2),(s,t_2),(t,t_2)\bigr).
\end{aligned}
\]
The companion mathematical-properties document uses
polynomial-major ordering
\[
\begin{aligned}
\mathcal O_{\mathrm{math}}={}&
\bigl((0,n),(0,t_1),(0,t_2),\\
& \;(s,n),(s,t_1),(s,t_2),\\
& \;(t,n),(t,t_1),(t,t_2)\bigr).
\end{aligned}
\]
The zero-based permutation between the two orderings is
\[
p = (0,3,6,1,4,7,2,5,8).
\]
Let $\Pi$ be the corresponding $9\times 9$ permutation matrix;
then $\trac_k^{\mathrm{math}} = \Pi\,\trac_k^{\mathrm{code}}$ and the
per-interface compliance block transforms as
$\DtN_k^{\mathrm{math}} = \Pi\,\DtN_k^{\mathrm{code}}\,\Pi^\top$.
Note that $\Pi = \Pi^\top = \Pi^{-1}$ (this particular permutation is
self-inverse).
The coefficient vector $\trac_k \in \R^9$ follows the code ordering
throughout this document; the mathematical properties document uses
the polynomial-major convention.

\textbf{Terminology.}  The variable $-\stress\normal$ appearing in the
traction Ansatz is the \emph{negative outward traction}, or
equivalently the inward-acting pressure-oriented traction.
We refer to it as ``inward traction'' for brevity; it is not a
distinct physical quantity but a sign-flipped pseudotraction.

% ======================================================================
\section{Affine traction Ansatz and code-compatible load sign}
% ======================================================================

\begin{assumption}[Local algebraic stability]\label{ass:local}
  For every pore $i$:
  \begin{enumerate}[label=(\alph*),nosep,leftmargin=*]
    \item $A_i$ is SPD.
    \item $B_i$ has full row rank (discrete inf--sup);
    \item the constant function $1$ belongs to the local pressure space.
  \end{enumerate}
\end{assumption}

For the standard conforming velocity space, condition~(a) follows, for
example, if $\pore{i}$ is connected and Lipschitz and its essential
no-slip boundary portion has positive surface measure.  The
Poincar\'e--Friedrichs inequality then excludes nonzero
constant-velocity null modes.

The unknown coefficients parameterize the \textbf{inward} traction:
\begin{equation}\label{eq:traction}
  -\stress(\vel_i,\pres_i)\,\normal_{i,k}
  = \sum_{(\alpha,\beta)\in\modeset}
    \lambda_{k,\alpha,\beta}\,\varphi_\alpha\,\bm{d}_{i,k,\beta}.
\end{equation}
The physical outward traction in the weak form is therefore
$\stress\normal_{i,k} = -\sum \lambda\varphi\bm{d}_{i,k,\beta}$.

The primitive load vector for mode $(\alpha,\beta)$ is
\begin{equation}\label{eq:load}
  \bigl[\loadvec_{i,k}^{(\alpha,\beta)}\bigr]_j
  = -\int_{\iface_k} \varphi_\alpha\;
    \bm{d}_{i,k,\beta}\cdot\bm{\phi}_j^{(i)}\,\mathrm{d}S.
\end{equation}
The leading minus sign is essential: it converts the inward-traction
Ansatz into the physical outward-traction boundary integral.
For the normal mode this gives the assembled direction $-\normal_{i,k}$;
for the tangential modes \eqref{eq:dir} gives $+\sigma_{i,k}\bm{t}_{k,\nu}$,
exactly the convention used in \texttt{affine\_face\_basis.py}.

At an inlet/outlet port only the constant normal mode is retained with
prescribed coefficient $\lambda_{i,\gamma}^k = p_b$.

% ======================================================================
\section{Local response and compliance map}
% ======================================================================

For each primitive mode solve
\begin{equation}\label{eq:response}
  \begin{bmatrix} A_i & B_i^\top \\ B_i & 0 \end{bmatrix}
  \begin{bmatrix} U_{i,k}^{(\alpha,\beta)} \\
                   X_{i,k}^{(\alpha,\beta)} \end{bmatrix}
  = \begin{bmatrix} \loadvec_{i,k}^{(\alpha,\beta)} \\ \zero \end{bmatrix}.
\end{equation}
Let $n_{u,i}$ denote the number of local velocity DOFs in $\pore{i}$.
Let $L_i \in \R^{n_{u,i}\times n_i}$ collect all $n_i = 9m_i^u + m_i^k$
primitive load columns.  Define
\begin{equation}\label{eq:H}
  M_i = B_i A_i^{-1} B_i^\top,\qquad
  H_i = A_i^{-1} - A_i^{-1}B_i^\top M_i^{-1} B_i A_i^{-1}.
\end{equation}
Then $U_i = H_i L_i \trac_i$ and the local compliance matrix is
\begin{equation}\label{eq:G}
  \DtN_i = L_i^\top H_i L_i \;\in\; \R^{n_i\times n_i},
\end{equation}
symmetric positive semidefinite, with
$\trac_i^\top \DtN_i \trac_i = U_i^\top A_i U_i \ge 0$.

For mode $(\alpha,\beta)$ on interface $\iface_k$, the \textbf{weighted
velocity moment} is
\begin{equation}\label{eq:moment}
  W_{i,k}^{(\alpha,\beta)}(U_i)
  = \int_{\iface_k} \varphi_\alpha\;
    U_{i,h}\cdot\bm{d}_{i,k,\beta}\,\mathrm{d}S.
\end{equation}
Only $W_{i,k}^{(0,n)}$ is the volumetric flow rate; the other eight are
weighted velocity moments (linear normal flux moments and tangential
velocity moments).  From the load convention,
$[\DtN_i\trac_i]_{k,\alpha,\beta} = -W_{i,k}^{(\alpha,\beta)}(U_i)$.
$\DtN_i$ is partitioned into $\DtN_i^{uu}$ (internal--internal) and
$\DtN_i^{uk}$ (internal--boundary).

% ======================================================================
\section{Global Schur complement assembly}
% ======================================================================

Continuity of all nine weighted velocity moments,
$W_{i,k}^{(\alpha,\beta)} + W_{j,k}^{(\alpha,\beta)} = 0$, together with
Boolean assembly gives
\begin{equation}\label{eq:schur}
  \Schur\,\trac = \bm{F},\qquad
  \Schur = \sum_{i=1}^N (R_i^u)^\top \DtN_i^{uu} R_i^u,\qquad
  \bm{F} = -\sum_{i=1}^N (R_i^u)^\top \DtN_i^{uk} R_i^k \trac^k.
\end{equation}
$\Schur \in \R^{9m\times 9m}$ is symmetric by construction.

\begin{remark}[Conditions for SPD]\label{rem:spd}
  $\Schur$ is positive definite under the following \emph{conditional}
  hypotheses (see the companion mathematical properties document for
  full statements and proofs):
  \begin{enumerate}[nosep,leftmargin=*]
    \item \textbf{Local stability:} Assumption~\ref{ass:local}.
    \item \textbf{Local kernel:} $\rank(H_i L_i) = n_i-1$ for every pore.
      This condition cannot be reduced to per-face linear independence
      of the nine modes; it additionally excludes cross-interface
      dependencies and dependencies introduced by the incompressible
      projection $H_i$.
    \item \textbf{Graph anchoring:} every connected component of the pore
      adjacency graph contains at least one inlet/outlet boundary port.
  \end{enumerate}
  \textbf{Floating basins} (pores without a solid-wall facet) violate
  Assumption~\ref{ass:local}(a) and are merged into neighbours before
  assembly.  This restores local coercivity but does \emph{not} supply
  condition~(3); the two issues are logically distinct.
\end{remark}

\medskip\noindent
\textbf{Constant-normal indicator vector.}
For each pore $i$, define $\one_{i,0n} \in \R^{n_i}$ by
\[
  [\one_{i,0n}]_\mu =
  \begin{cases}
    1, & \text{if mode $\mu$ is the constant-normal mode $(0,n)$
         on some port of $\pore{i}$},\\
    0, & \text{otherwise}.
  \end{cases}
\]
This vector has entry $1$ on every port's $(0,n)$ slot (internal and
boundary alike) and $0$ on all other slots.  In code, the diagnostic
mask must be constructed as this specific vector; using
\texttt{np.ones(}$n_i$\texttt{)} would produce an incorrect all-ones
vector that does \emph{not} represent the constant-pressure kernel.

% ======================================================================
\section{Algorithm}
% ======================================================================

\begin{center}
\fbox{%
\begin{minipage}{0.92\textwidth}
\smallskip
\noindent\textbf{Algorithm 1:} Affine-DDPNM solver

\smallskip\noindent\textbf{Input:}
Subdomain meshes $\{\pore{i}\}$, FE pairs $(\mathbf{V}_{i,h},Q_{i,h})$,
prescribed boundary coefficients $\trac^k$.

\noindent\textbf{Output:}
Unknown interface coefficients $\trac \in \R^{9m}$, local fields
$\{(U_i,P_i)\}$, weighted velocity moments, and residuals.

\smallskip
\noindent\emph{Offline library} --- for each pore $i=1,\dots,N$ (in parallel):
\begin{enumerate}[nosep,leftmargin=18pt]
  \item Assemble $A_i$, $B_i$.
  \item For each internal interface $\gamma \in \mathcal{U}_i$ and each
    mode $(\alpha,\beta) \in \modeset$:
    assemble $\bm{f}_{i,\gamma}^{(\alpha,\beta)}$ via~\eqref{eq:load},
    solve~\eqref{eq:response} for
    $(U_{i,\gamma}^{(\alpha,\beta)}, X_{i,\gamma}^{(\alpha,\beta)})$.
  \item For each boundary port $\gamma \in \mathcal{K}_i$:
    assemble and solve the single $(0,n)$ mode.
  \item Form $L_i$, compute $\DtN_i = L_i^\top H_i L_i$ via~\eqref{eq:G}.
    As a diagnostic, compute the singular values of $H_i L_i$ and confirm
    one numerically zero singular value (within a stated tolerance) with
    its right singular vector aligned with $\one_{i,0n}$; equivalently,
    inspect the eigenvalues of $\DtN_i$ and confirm one numerically zero
    eigenvalue with the corresponding eigenvector aligned with
    $\one_{i,0n}$.  (This is the local-kernel condition in
    Remark~\ref{rem:spd}, item~2.  \textbf{It is an assumption, not a
    theorem:} it is not proved for arbitrary pore geometries or meshes
    and must be verified numerically for each local library.)
  \item Accumulate global Schur system from the raw assembled blocks
    (before any explicit symmetrization):
    $\Schur \mathrel{+}= (R_i^u)^\top \DtN_i^{uu} R_i^u$,
    $\bm{F} \mathrel{-}= (R_i^u)^\top \DtN_i^{uk} R_i^k \trac^k$.
\end{enumerate}

\noindent\emph{Global solve}:
Check the symmetry defect $\|\Schur - \Schur^\top\| / \|\Schur\|$ of the
raw assembled matrix; only then apply $\Schur \leftarrow
\frac12(\Schur+\Schur^\top)$ if needed.  Verify positive definiteness
(Cholesky factorization or smallest eigenvalue).  Solve
$\Schur\,\trac = \bm{F}$ by sparse direct solver.

\noindent\emph{Reconstruction} --- for each pore $i=1,\dots,N$ (in parallel):
\[
  \textstyle
  U_i = \sum_{\gamma\in\mathcal{U}_i}\sum_{\alpha,\beta}
    \lambda_{i,\gamma,\alpha,\beta}\,U_{i,\gamma}^{(\alpha,\beta)}
    + \sum_{\gamma\in\mathcal{K}_i} \lambda_{i,\gamma}^k\,
    U_{i,\gamma}^{(0,n)},
\]
and analogously for $P_i$ with the pressure responses $X_{i,\gamma}^{(\cdot)}$.

\noindent\emph{Verification} (compute independently):
\begin{align*}
  r_{B,i} &= \|B_i U_i\| \,/\, (\|U_i\| + \epsilon_{\text{mach}})
    \qquad \text{(local discrete incompressibility)}, \\[2pt]
  r_{\text{port},i} &= \Bigl|\sum_{\gamma\subset\partial\pore{i}}
    W_{i,\gamma}^{(0,n)}\Bigr|
    \qquad \text{(per-pore mass balance)}, \\[2pt]
  r_{\text{flux}} &= \max_{\gamma}\; \bigl|
    W_{i,\gamma}^{(0,n)} + W_{j,\gamma}^{(0,n)}\bigr|
    \qquad \text{(interface flux continuity)}, \\[2pt]
  r_{\text{mom}} &= \max_{\gamma,\alpha,\beta}\; \bigl|
    W_{i,\gamma}^{(\alpha,\beta)} + W_{j,\gamma}^{(\alpha,\beta)}\bigr|
    \qquad \text{(all nine moment continuity)}.
\end{align*}
Here $W_{i,\gamma}^{(\alpha,\beta)}$ denotes the weighted velocity moment
(equation~\eqref{eq:moment}); only $W_{i,\gamma}^{(0,n)}$ is the
volumetric flow rate.  The implementation reports
$r_{\text{mom}}$ as \texttt{max\_mass\_residual} (computed from the full
moment-residual vector and covering all nine modes); the naming is
legacy and refers to the same quantity.
\smallskip
\end{minipage}}
\end{center}

% ======================================================================
\section{Relation to classic DD-PNM and scope of enrichment}
% ======================================================================

Let $E_{\text{P0}} : \R^m \to \R^{9m}$ embed each classic coefficient
into the constant-normal slot.  The classic DD-PNM system is recovered
by restricting \textbf{both trial and test spaces}:
\[
  \Schur_{\text{P0}} = E_{\text{P0}}^\top \Schur E_{\text{P0}},\qquad
  \bm{F}_{\text{P0}} = E_{\text{P0}}^\top \bm{F}.
\]
An affine solution of
$\Schur\,\trac^{\mathrm{aff}} = \bm{F}$
generally does \emph{not} have the classic solution as its
constant-normal subvector, because the eight additional test-space
constraints alter the Galerkin projection.  Under the SPD hypotheses
of Remark~\ref{rem:spd} this affine solution is unique.

The valid enrichment statement is conditional and reference-dependent.
If a well-defined full interface-traction Schur system
$\widehat{\Schur}\,\bm{\xi} = \widehat{\bm{F}}$ is introduced with the
same local pressure spaces and sign conventions (the broken-pressure
C-DD system of the companion mathematical properties document), then
the affine system is its Galerkin restriction and the affine reduction
error is no larger than the classic reduction error in the
$\widehat{\Schur}$-energy seminorm.  By the corresponding
velocity-energy identity, the same monotonicity holds for the
broken-$H^1$ velocity reduction error measured relative to the full
C-DD solution.  This argument does \emph{not} imply monotonic
improvement of velocity $L^2$, pressure $L^2$, or outlet-flux errors,
nor of errors measured relative to the exact PDE solution.

No theorem in the present analysis establishes equality between that
broken-pressure traction system and the monolithic continuous-pressure
Taylor--Hood system.  The ``Exact FE-trace Schur'' computation is a
primal Schur reduction of the monolithic system.  For its $O(10^{-12})$
agreement with the direct monolithic solve to hold, the two must operate
on the \textbf{same assembled matrix}---identical pressure stabilisation,
pressure constraints (pinning / zero-mean), boundary conditions, and
gauge treatment.  If the monolithic reference uses $\varepsilon=0$ and
the Schur construction uses $\varepsilon=10^{-10}$, the matrices differ
and the comparison is not an exact equivalence of a single discrete
system.

% ======================================================================
\section{Computational cost}
% ======================================================================

A pore with $m_i^u$ internal ports and $m_i^k$ boundary ports requires
$n_i = 9m_i^u + m_i^k$ primitive solves.  Offline cost grows approximately
$9\times$ relative to the one-mode-per-interface classic method.
The global unknown count is $9m$; sparse factorization cost depends on
graph fill and cannot be inferred from the unknown-count factor alone.
For multi-boundary-condition studies the offline library is built once.

\begin{thebibliography}{9}
\bibitem{sun2025ddpnm}
S.~Sun, Z.~Wang, L.~Zhang, J.~Zhao.
\newblock A domain decomposition approach to pore-network modeling of porous
  media flow.
\newblock \textit{arXiv:2510.13429v2}, 2026.
\end{thebibliography}

\end{document}
